%%%%%
%
\section{Conversion des coordonn�es d'un vecteur 2D :
de polaire g\'eographique vers cart�sien}

\itemstyle{Op\'erateur~:} {\sf polgeo2rect}
\index{polgeo2rect}

%%%%%OBJET%%%%
\itemstyle{Objet~:}

Calcule les coordonn�es cart�siennes d'un vecteur 2D 
initialement exprim� dans 
un syst�me polaire g�ographique (angle compt� positif
dans le sens indirect � partir du nord g�ographique).


%%%%%ARGUMENTS%%%%
\bigskip
%%%%%ENTREES%%%%
\begin{tabarg}
\multicolumn{3}{l}{\titretab {Entr\'ees}}\\
\hline
$M$ & Vecteur, \texttt{float} & Module (any)\\
$alpha$ & Vecteur, \texttt{float} & Angle (degre ou radian)\\
$unite$ & Coeff., \texttt{String  } & Unit\'e de l'angle (typiquement \texttt{\$Arg2.unit})\\
\end{tabarg}

\medskip
%%%%%SORTIES%%%%
\begin{tabarg}
\multicolumn{3}{l}{\titretab {Sorties}}\\
\hline
$Comp_{Est}$ & Vecteur, \texttt{float} & Composante zonale (m\^eme unit\'e que $M$)\\
$Comp_{Nord}$ & Vecteur, \texttt{float} & Composante m\'eridienne (m\^eme unit\'e que $M$)\\
\end{tabarg}


%%%%%FORMULE%%%%
\itemstyle{Formule}

%%polgeo2rect
\begin{align*}
  Comp_{Est}  &=  M . \sin (alpha) \\
  Comp_{Nord} &=  M . \cos (alpha)
\end{align*}



%%%%%AUTEUR%%%%
\itemstyle{Auteur~:}
CNRM/GMEI/TRAMM

%%%%%REFERENCE%%%%
\itemstyle{R\'ef\'erences~:}

